\section{低维模型的拟合与前瞻预测}
\label{sec:prediction}
\subsection{资源模型与积累模型}
状态干预表明，资源变量在M1的响应递减中起重要作用。若研究目标仅为aBN1窗口计数，是否仍需保留全脑状态，是另一个需要检验的问题。为此，构造资源模型S与积累模型I，近似描述M1的输入至计数关系。两者均为独立的低维代理，不改变全脑模型。

令$f$为经过输入接口转换后的平均感觉发放率，$u=f/150$，$P$为刺激持续时间（秒），横线表示刺激期间的状态均值，$[a]_+=\max(0,a)$。刺激期间有
\begin{align}
\dot x&=(1-x)/\tau-Ufx,&\widehat R_S&=\frac{P}{0.2}[a_Su\bar x-b_S]_+,\label{eq:S}\\
\dot z&=(u-z)/\tau,&\widehat R_I&=\frac{P}{0.2}[a_Iu-b_I\bar z]_+.\label{eq:I}
\end{align}
静默时，S的资源向1恢复，I的积累状态向0衰减。两个模型均固定$\tau=2$s，S另固定$U=0.01$，各拟合两个非负读出系数。常输入阶段以解析式推进状态，计算刺激内均值后代入读出函数。模型不显式表示细胞间连接、单个脉冲时间或刺激后100ms内的活动。

持续时间缩放项$P/0.2$在前瞻实验前确定，用于从200ms刺激外推到其他宽度。S的状态方程与全脑M1的资源机制相近，因而具有有利的先验。I则采用另一种受限的状态与读出形式；比较结果仅适用于这两个具体模型。

\subsection{参数估计与回顾性检查}
参数拟合仅使用首轮M1短间隔条件的60个训练窗口，其余80个M1窗口用于跨条件检查。两部分均属于首轮280个窗口，未增加观测数量。由于这80个窗口的结果此前已经读取，本文将其列为回顾性检查。

\begin{table}[!htbp]\centering\small
\caption{低维模型的拟合系数及旧数据误差。RMSE单位为脉冲。两个模型各估计两个系数，其状态方程与读出函数不同。}
\begin{tabular}{lrrrr}\toprule
模型 & $a$ & $b$ & 拟合RMSE & 回顾RMSE\\\midrule
\inputrows{data/fit.tex}\bottomrule
\end{tabular}\end{table}

S的拟合RMSE和回顾性RMSE分别为0.490和0.774脉冲；I分别为0.732和0.939脉冲。两者都能较好地近似旧数据。尽管状态干预已确定资源变量对M1响应的作用，积累模型仍可拟合相近的计数曲线，说明仅依靠这一输出序列难以唯一识别内部动力学。新条件下的预测比较因此尤为重要。

\subsection{判别刺激的选择与预测保存}
候选刺激由100/150/200Hz输入率、100/200/400ms宽度和刺激后200/800ms静默组成，共18组。以两个模型预测之间的RMSE衡量分歧，先选择分歧最大的条件，再选取频率或宽度不同的下一条件，另设置刺激形状不变、仅更换种子的对照。

\begin{table}[!htbp]\centering\small
\caption{前瞻实验的刺激条件。$P$为刺激宽度，$W$为计数窗口，$T$为起始间隔。各条件单独评分，以免混合不同长度窗口的计数。}
\begin{tabular}{lrrrrr}\toprule
条件 & 命令Hz & $P$(ms) & $W$(ms) & $T$(s) & 重复数\\\midrule
新种子对照 & 150 & 200 & 300 & 0.5 & 12\\
200Hz长刺激 & 200 & 400 & 500 & 1.2 & 8\\
150Hz长刺激 & 150 & 400 & 500 & 1.2 & 8\\\bottomrule
\end{tabular}\end{table}

前瞻实验采用2701--2703，训练后分别从同一末态等待1s和5s探测。模型系数、刺激选择和102个M1窗口的逐点预测均在新全脑仿真产生输出之前保存，并记录哈希。另运行9个M0初态单次刺激窗口作为传递对照，不参与拟合、条件选择或预测评分；这些窗口不构成完整的M0重复训练对照。

预设准确度标准为每个条件的平均绝对误差（MAE）不超过1.5脉冲，同时比较两个判别条件下的MAE与均方根误差（RMSE）。由于实验主动选择了预测分歧较大的刺激，误差比较主要反映两模型在判别条件下的表现，不用于估计自然刺激分布上的平均性能。

\subsection{前瞻预测结果}
\begin{table}[!htbp]\centering\small
\caption{前瞻预测误差。各条件包含三个输入实现的训练及恢复窗口，达标指MAE$\leq1.5$脉冲。模型未使用这些输出重新拟合。}
\label{tab:forecast}
\begin{tabular}{llrrrr}\toprule
条件 & 模型 & 窗口数 & MAE & RMSE & 达标\\\midrule
\inputrows{data/forecasts.tex}\bottomrule
\end{tabular}\end{table}

S在三个条件中均满足准确度标准，I仅在新种子对照中满足。在两个判别条件下，S的MAE和RMSE均低于I；合并102个窗口，总体RMSE分别为1.036与3.019脉冲。所有输入实现均纳入评分。

\newcommand{\ForecastPanel}[2]{\begin{tikzpicture}\begin{axis}[width=.94\linewidth,height=5cm,title={#2},xlabel={刺激序号},ylabel={响应计数},xmin=1,ymin=0,grid=major,grid style={black!10},tick label style={font=\small},title style={font=\small}]
\addplot[black,thick,mark=*] table[x=trial,y=aBN1]{data/forecast_#1.csv};
\addplot[Blue,thick,dashed,mark=triangle*] table[x=trial,y=resource]{data/forecast_#1.csv};
\addplot[Orange,thick,dashdotted,mark=square*] table[x=trial,y=accumulator]{data/forecast_#1.csv};
\end{axis}\end{tikzpicture}}
\begin{figure}[!htbp]\centering
\begin{minipage}{.48\linewidth}\ForecastPanel{1}{200Hz，400ms}\end{minipage}\hfill
\begin{minipage}{.48\linewidth}\ForecastPanel{2}{150Hz，400ms}\end{minipage}
\par\smallskip\small 黑色：全脑M1\quad\textcolor{Blue}{蓝色：资源S}\quad\textcolor{Orange}{橙色：积累I}
\caption{判别条件下的实测响应与预测。曲线为三个输入实现的训练响应均值；完整评分还包括表\ref{tab:forecast}所列的恢复探测窗口。}
\end{figure}

在所测试条件下，具有资源动力学先验的低维模型较准确地预测了M1的平均响应。然而，这一比较未涵盖其他时间尺度、非线性读出或抑制性可塑性模型，故不足以排除这些机制。

\subsection{强刺激下的局部响应回升}
200Hz长刺激的首个输入实现出现4$\to$6脉冲的相邻响应回升。观察到这一变化后，对全部九条新序列补充统计最大相邻增量。此项分析属于事后探索，未参与原条件选择、参数估计或准确度判定。

三个200Hz实现均出现局部回升，另外两个条件未见上升。全部九条序列保存在分析表中；下表列出强刺激条件下的最大回升及对应资源变化。

\begin{table}[!htbp]\centering\small
\caption{200Hz长刺激下的相邻响应回升。资源变化指M1刺激前边资源均值之差，S预测变化指对应相邻刺激的预测计数之差。}
\begin{tabular}{rrrrr}\toprule
种子 & 刺激序号 & 实测计数 & 资源均值变化 & S预测变化\\\midrule
\inputrows{data/rebounds.tex}\bottomrule
\end{tabular}\end{table}

回升发生时，平均资源仍在下降，S也预测响应略降。两种低维模型在固定重复刺激下均只产生下降或平台，未能准确再现这些局部变化。因此，平均资源与简单单调读出不足以描述全部计数波动。潜在因素包括输入时序、细胞间资源分配、兴奋与抑制的相互作用以及阈值重置的非线性，但当前未通过选择性干预区分这些解释。

S在总体误差上的表现与其局部预测偏差可以同时成立。现有结果说明其适用于部分窗口计数预测，尚未说明精确发放时序、细胞干预后果或完整连接结构的必要性。
\FloatBarrier
